\setcounter{section}{2}

  \zwischenueberschrift{Übungsaufgaben}

\inputaufgabegibtloesung
{}
{

Es seien
\maabbdisp {f,g} {\R} { \R^n
} {}
\definitionsverweis {differenzierbare Kurven}{}{.}
Berechne die
\definitionsverweis {Ableitung}{}{}
der Funktion
\maabbeledisp {} {\R} {\R
} { t } { \left\langle f(t) , g(t)  \right\rangle
} {.}
Formuliere das Ergebnis mit dem Skalarprodukt.

}
{} {}

\inputaufgabegibtloesung
{}
{

Es sei
\maabb {h} { I } { V
} {}
eine zweimal differenzierbare Kurve in einem
\definitionsverweis {euklidischen Vektorraum}{}{}
$V$. Zeige, dass bei

\mavergleichskette
{\vergleichskette
{h'(t)
}
{ \neq }{ 0
}
{  }{
}
{    }{
}
{   }{
}
}
{}{}{}
die Gleichheit

\mavergleichskettedisp
{\vergleichskette
{ \Vert { h'(t)} \Vert'
}
{ =} { { \frac{   \left\langle h'(t) , h^{\prime \prime} (t)  \right\rangle  }{  \left\langle h'(t) , h'(t)  \right\rangle  } } \Vert { h'(t)} \Vert
}
{ } {
}
{ } {
}
{ } {
}
}
{}{}{}
gilt.

}
{} {}

Eine
\definitionsverweis {differenzierbare Kurve}{}{}
\maabbeledisp {f} { [a,b] } {\R^n
} { t } {f(t) = \left( f_1(t)  , \,   \ldots  , \,  f_n(t)                 \right)
} {,}
heißt
\definitionswort {bogenparametrisiert}{,}
wenn

\mavergleichskette
{\vergleichskette
{ f_1'(t)^2 + \cdots +  f_n'(t)^2
}
{ = }{ 1
}
{  }{
}
{    }{
}
{   }{
}
}
{}{}{}
für alle $t$ gilt.

\inputaufgabe
{}
{

Es sei
\maabb {\gamma} { I } {\R^n
} {}
eine zweimal
\definitionsverweis {stetig differenzierbare}{}{}
\definitionsverweis {bogenparametrisierte}{}{}
Kurve. Zeige, dass
\mathkor {} {\gamma'(t)} {und} {\gamma^{\prime \prime}(t)} {}
aufeinander senkrecht stehen.

}
{} {}

\inputaufgabe
{}
{

Es sei

\mavergleichskette
{\vergleichskette
{ P
}
{ \in }{ \R^n
}
{  }{
}
{    }{
}
{   }{
}
}
{}{}{}
ein Punkt und sei

\mavergleichskette
{\vergleichskette
{ I
}
{ = }{ {]{-1},1[}
}
{  }{
}
{    }{
}
{   }{
}
}
{}{}{.}
Wir betrachten die Menge

\mavergleichskettedisp
{\vergleichskette
{ M
}
{ =} { \left\{ f:I \rightarrow \R^n   \mid  f \text{ differenzierbar} , \,  f(0) = P       \right\}
}
{ } {
}
{ } {
}
{ } {
}
}
{}{}{.}
Wir nennen zwei
\definitionsverweis {Kurven}{}{}

\mavergleichskette
{\vergleichskette
{ f,g
}
{ \in }{ M
}
{  }{
}
{    }{
}
{   }{
}
}
{}{}{}
\stichwort {tangential äquivalent} {,} wenn

\mavergleichskettedisp
{\vergleichskette
{ f'(0)
}
{ =} { g'(0)
}
{ } {
}
{ } {
}
{ } {
}
}
{}{}{}
ist.

a) Zeige, dass dies eine
\definitionsverweis {Äquivalenzrelation}{}{}
ist.

b) Finde den einfachsten Vertreter für die
\definitionsverweis {Äquivalenzklassen}{}{.}

c) Man gebe für jede Klasse einen weiteren Vertreter an.

d) Beschreibe die Menge der Äquivalenzklassen
\zusatzklammer {also die
\definitionsverweis {Quotientenmenge}{}{}} {} {.}

}
{} {}

\inputaufgabe
{}
{

Man gebe ein Beispiel für eine
\zusatzklammer {überall reguläre} {} {}
\definitionsverweis {differenzierbare Hyperfläche}{}{,}
die nicht
\definitionsverweis {zusammenhängend}{}{}
ist.

}
{} {}

\inputaufgabe
{}
{

Es sei

\mavergleichskette
{\vergleichskette
{ V
}
{ \subseteq }{ \R^{n-1}
}
{  }{
}
{    }{
}
{   }{
}
}
{}{}{}
eine
\definitionsverweis {offene Menge}{}{}
und sei
\maabbdisp {f} { V } {\R
} {}
eine
\definitionsverweis {stetig differenzierbare}{}{}
Funktion, wir betrachten den
\definitionsverweis {Graphen}{}{}
$Y$ zu $f$ als differenzierbare Hyperfläche in $\R^n$ im Sinne von
[[Funktion/n-1 Variablen/Graph als Hyperfläche/Tangentialraum/Beispiel|Beispiel 1.2]].
Bestimme die
\definitionsverweis {Einheitsnormalenfelder}{}{}
in diesem Fall.

}
{} {}

\inputaufgabegibtloesung
{}
{

Es seien
\mathkor {} {r} {und} {R} {}
positive reelle Zahlen mit

\mavergleichskette
{\vergleichskette
{ 0
}
{ < }{ r
}
{ < }{ R
}
{    }{
}
{   }{
}
}
{}{}{.}
Bestimme ein
\definitionsverweis {Einheitsnormalenfeld}{}{}
für den
\zusatzklammer {eingebetteten} {} {}
Torus

\mavergleichskettedisp
{\vergleichskette
{ T
}
{ =} { \left\{ (x,y,z) \in \R^3  \mid   \left(  \sqrt{x^2+y^2} -R  \right)^2 +z^2  = r^2         \right\}
}
{ } {
}
{ } {
}
{ } {
}
}
{}{}{.}

}
{} {}

\inputaufgabe
{}
{

Es sei
\maabbdisp {f} { I } {\R_+
} {}
eine
\definitionsverweis {differenzierbare Funktion}{}{}
und sei $Y$ die zugehörige Rotationsfläche zum Graphen von $f$ als differenzierbare Hyperfläche in $\R^n$ im Sinne von
[[Rotationsfläche/Differenzierbare positive Kurve/Graph/Tangentialraum/Fakt|Lemma 2.1]].
Bestimme die
\definitionsverweis {Einheitsnormalenfelder}{}{}
in diesem Fall.

}
{} {}

\inputaufgabe
{}
{

Zeige, dass die
\definitionsverweis {Gauß-Abbildung}{}{}
auf einer Hyperebene konstant ist. Gilt davon auch die Umkehrung?

}
{} {}

\inputaufgabe
{}
{

Zeige, dass die
\definitionsverweis {Gauß-Abbildung}{}{}
auf einer Sphäre im $\R^n$ mit dem Ursprung als Mittelpunkt durch eine Streckung des $\R^n$ induziert wird.

}
{} {}

\inputaufgabe
{}
{

Es sei

\mavergleichskette
{\vergleichskette
{ Y
}
{ \subseteq }{ \R^n
}
{  }{
}
{    }{
}
{   }{
}
}
{}{}{}
eine
\definitionsverweis {differenzierbare Hyperfläche}{}{}
mit dem
\definitionsverweis {Einheitsnormalenfeld}{}{}
$N$ und dem entgegengesetzten Einheitsnormalenfeld $-N$. Zeige, dass sich die zugehörigen
\definitionsverweis {Gauß-Abbildungen}{}{}
um die
\definitionsverweis {Antipodenabbildung}{}{}
\maabbeledisp {} {S^{n-1}} {S^{n-1}
} { P } { -P
} {m}
unterscheiden.

}
{} {}

\inputaufgabegibtloesung
{}
{

Es seien

\mavergleichskette
{\vergleichskette
{ \alpha,\beta
}
{ \in }{ \R_+
}
{  }{
}
{    }{
}
{   }{
}
}
{}{}{}
und sei

\mavergleichskettedisp
{\vergleichskette
{ Y
}
{ =} { \left\{ (x,y)\in\R^2  \mid   \alpha x^2+\beta y^2 = 1        \right\}
}
{ } {
}
{ } {
}
{ } {
}
}
{}{}{.}
\aufzaehlungzwei {Bestimme die
\definitionsverweis {Gauß-Abbildung}{}{}
zu $Y$.
} { Man gebe zur Gauß-Abbildung zu $Y$ explizit eine
\definitionsverweis {Umkehrabbildung}{}{}
an.
}

}
{} {}

\inputaufgabegibtloesung
{}
{

Es seien

\mavergleichskette
{\vergleichskette
{ \alpha,\beta
}
{ \in }{ \R_+
}
{  }{
}
{    }{
}
{   }{
}
}
{}{}{}
und sei

\mavergleichskettedisp
{\vergleichskette
{ Y
}
{ =} { \left\{ (x,y)\in\R^2  \mid   \alpha x^2+\beta y^2 = 1        \right\}
}
{ } {
}
{ } {
}
{ } {
}
}
{}{}{.}
Zeige, dass die
\definitionsverweis {Gauß-Abbildung}{}{}
zu $Y$ bijektiv ist.

}
{} {}

\inputaufgabe
{}
{

Man gebe ein Beispiel für eine
\definitionsverweis {differenzierbare Hyperfläche}{}{}

\mavergleichskette
{\vergleichskette
{ Y
}
{ \subseteq }{ \R^2
}
{  }{
}
{    }{
}
{   }{
}
}
{}{}{}
derart, dass die
\definitionsverweis {Gauß-Abbildung}{}{}
surjektiv ist und dass es Punkte auf $S^1$ gibt, die unendlich oft angenommen werden.

}
{} {}

\inputaufgabe
{}
{

Es sei

\mavergleichskette
{\vergleichskette
{ W
}
{ \subseteq }{ \R^n
}
{  }{
}
{    }{
}
{   }{
}
}
{}{}{}
\definitionsverweis {offen}{}{,}
\maabb {h} { W } { \R
} {}
eine
\definitionsverweis {stetig differenzierbare Funktion}{}{}
und

\mavergleichskette
{\vergleichskette
{ Y
}
{ = }{ h^{-1}(c)
}
{  }{
}
{    }{
}
{   }{
}
}
{}{}{}
die
\definitionsverweis {Faser}{}{}
zu

\mavergleichskette
{\vergleichskette
{ c
}
{ \in }{ \R
}
{  }{
}
{    }{
}
{   }{
}
}
{}{}{,}
wobei $h$ in jedem Punkt von $Y$
\definitionsverweis {regulär}{}{}
sei. Es sei
\maabbdisp {L} {\R^n } { \R^n
} {}
eine
\definitionsverweis {lineare Isometrie}{}{,}

\mavergleichskette
{\vergleichskette
{ W'
}
{ = }{ L^{-1}(W)
}
{  }{
}
{    }{
}
{   }{
}
}
{}{}{}
und

\mavergleichskettedisp
{\vergleichskette
{ Z
}
{ =} { L^{-1}(Y)
}
{ =} { (h \circ L)^{-1}(c)
}
{ } {
}
{ } {
}
}
{}{}{.}
Zeige, dass sich die Konzepte geodätische Kurve, Normalenfeld, Einheitsnormalenfeld, Gauß-Abbildung auf $Y$ und auf $Z$ entsprechen
\zusatzklammer {also mithilfe von $L$ ineinander überführt werden können} {} {.}

}
{} {}

  \zwischenueberschrift{Aufgaben zum Abgeben}

\inputaufgabe
{3}
{

Bestimme das Bild unter der
\definitionsverweis {Gauß-Abbildung}{}{}
zur Standardparabel mit der nach oben gerichteten Orientierung.

}
{} {}

\inputaufgabe
{4}
{

Es sei

\mavergleichskette
{\vergleichskette
{ W
}
{ = }{ \R^3 \setminus \{ (0,0,0) \}
}
{  }{
}
{    }{
}
{   }{
}
}
{}{}{}
und sei

\mavergleichskette
{\vergleichskette
{ Y
}
{ \subseteq }{ W
}
{  }{
}
{    }{
}
{   }{
}
}
{}{}{}
die
\definitionsverweis {Nullstellenmenge}{}{}
zu

\mavergleichskettedisp
{\vergleichskette
{ h(x,y,z)
}
{ =} { x^2+y^2-z^2
}
{ } {
}
{ } {
}
{ } {
}
}
{}{}{.}
Bestimme das Bild von $Y$
\zusatzklammer {mit der durch den Gradienten von $h$ gegebenen Orientierung} {} {}
unter der
\definitionsverweis {Gauß-Abbildung}{}{.}

}
{} {}

\inputaufgabe
{5}
{

Es seien

\mavergleichskette
{\vergleichskette
{ \alpha,\beta, \gamma
}
{ \in }{ \R_+
}
{  }{
}
{    }{
}
{   }{
}
}
{}{}{}
und sei

\mavergleichskettedisp
{\vergleichskette
{ Y
}
{ =} { \left\{ (x,y,z) \in \R^3  \mid   \alpha x^2+ \beta y^2 + \gamma z^2 = 1        \right\}
}
{ } {
}
{ } {
}
{ } {
}
}
{}{}{.}
Zeige, dass die
\definitionsverweis {Gauß-Abbildung}{}{}
zu $Y$ bijektiv ist.

}
{} {}

\inputaufgabe
{3}
{

Es sei

\mavergleichskette
{\vergleichskette
{ W
}
{ \subseteq }{ \R^n
}
{  }{
}
{    }{
}
{   }{
}
}
{}{}{}
\definitionsverweis {offen}{}{,}
\maabb {h} { W } { \R
} {}
eine
\definitionsverweis {stetig differenzierbare Funktion}{}{}
und

\mavergleichskette
{\vergleichskette
{ Y
}
{ = }{ h^{-1}(c)
}
{  }{
}
{    }{
}
{   }{
}
}
{}{}{}
die
\definitionsverweis {Faser}{}{}
zu

\mavergleichskette
{\vergleichskette
{ c
}
{ \in }{ \R
}
{  }{
}
{    }{
}
{   }{
}
}
{}{}{,}
wobei $h$ in jedem Punkt von $Y$
\definitionsverweis {regulär}{}{}
sei. Es sei
\maabbdisp {L} {\R^n } { \R^n
} {}
eine
\definitionsverweis {lineare Isomorphie}{}{,}

\mavergleichskette
{\vergleichskette
{ W'
}
{ = }{ L^{-1}(W)
}
{  }{
}
{    }{
}
{   }{
}
}
{}{}{}
und

\mavergleichskettedisp
{\vergleichskette
{ Z
}
{ =} { L^{-1}(Y)
}
{ =} { (h \circ L)^{-1}(c)
}
{ } {
}
{ } {
}
}
{}{}{.}
Zeige, dass sich die Konzepte geodätische Kurve, Normalenfeld, Einheitsnormalenfeld, Gauß-Abbildung auf $Y$ und auf $Z$ im Allgemeinen nicht entsprechen
\zusatzklammer {im Gegensatz zu
[[Isometrische Isomorphie/Differenzierbare Hyperfläche/Geodätische, Normalenfeld, Gauß-Abbildung/Aufgabe|Aufgabe 2.14]],
wo eine Isometrie vorliegt} {} {.}

}
{} {}

 [[Kategorie:Latexseite]]
